\documentclass[11pt,answers]{exam}
%\usepackage{currfile}

\usepackage{amsmath,amssymb,amsthm}
\DeclareFontEncoding{T1}
\fontencoding{T1}
\usepackage[T2A]{fontenc}
\usepackage[utf8]{inputenc}
\usepackage[russian]{babel}
\usepackage{graphics,graphicx}
\usepackage{indentfirst}

\usepackage{pythontex}

\usepackage{verbatim}
\usepackage{fancyvrb}
\usepackage{booktabs}
\usepackage{multirow}

\usepackage{listings}
\usepackage{color}

\definecolor{gray}{gray}{0.5}

\lstset{basicstyle=\ttfamily,
    language=python,    
    keepspaces=true,
    extendedchars=\true,
    identifierstyle={ \bf \color{black} },
    showstringspaces=false,
    numbers=left,%
    numbersep=7pt, 
    emph=[1]{case,switch,otherwise,nonlocal,as,yield,zip,print},    
    frame=single,    
    xleftmargin=0.3cm,
    frame=l,framesep=4pt,framerule=0.5pt,
    numberstyle={\scriptsize \color{gray}}
}


\pagestyle{headandfoot}
\runningheadrule
\firstpageheader{ИнтМатПак}{MATLAB 4.\jobname}{24-31 марта 2026}
\runningheader{ИнтМатПак}
{MATLAB 4.\jobname, Стр. \thepage\ из \numpages}
{24-31 марта 2026}
\firstpagefooter{}{}{}
\runningfooter{}{}{}
\pointname{} 
\boxedpoints
%\pointsinmargin
\pointsinrightmargin

\begin{document}

\vspace{0.1in}
\makebox[\linewidth]{Ф. И. О. \enspace \hrulefill \enspace  группа: \enspace\hrulefill}
\vspace{0.1in}

\begin{questions}

\question[1] Найдите решение системы линейных уравнений и запишите значение $x_1$:

\begin{pylabcode} 
from numpy import *
import scipy.integrate as integrate
from scipy.optimize import fsolve

n = 4 
A = (random.rand(n,n)*10+1).round()
res1 = (random.rand(n,1)*10).round()
B = dot(A,res1)
sA = '\\begin{align}\n'
for i in range(A.shape[0]):
    for j in range(A.shape[1]):
        sA = sA + '{:2.0f}'.format(A[i,j]) + '\cdot x_' + str(j) + '+'
    else:         
        sA = sA[0:-1]
        sA = sA + ' = ' + '{:3.0f}'.format(B[i,0]) + '\\\\'
else:
    sA = sA[0:-2]
    sA = sA + '\n\\end{align}'    
\end{pylabcode} 
\pylab{sA}

\answerline[$x_1 = $\color{white}{\pylab{res1[1][0]}}] 
 
\question[2]
\begin{pylabcode} 
x = array(range(0,5)) 
y = 10*sin(x*0.5*pi/max(x))+random.rand(x.shape[0]).round(2)
z = polyfit(x, y, 3)
p = poly1d(z)
res2 = p(2.6)
\end{pylabcode}
Задана табличная функция:
$$
x =  
\begin{pylabcode} 
print(array2string(x.flatten(),separator=', ',precision=1)) 
\end{pylabcode}
, \quad y = 
\begin{pylabcode} 
print(array2string(y.flatten(),separator=', ',precision=2))
\end{pylabcode}
$$
Напишите код и вычислите значение функции при x=2.6, аппроксимировав табличную функцию многочленом третьей степени. 
\fillwithdottedlines{3cm}
\answerline[$y(2.6) =$ \color{white}{\pylab{'{:4.2f}'.format(res2)}}] 


\question[2] Найдите ближайший к нулю корень уравнения:
\begin{pylabcode} 
a = int((random.rand(1,1)*10+2)[0])
b = int((random.rand(1,1)*3+3)[0])
res3 = fsolve(lambda x: cos(4*x)-10*x, 0)
\end{pylabcode}
$$
    \cos \pylab{b} x - \pylab{a} x = 0
$$
\answerline[$x=$ \color{white}{\pylab{'{:4.2f}'.format(res3[0])}}] 

\question[2] Задана табличная функция:
\begin{pylabcode} 
x = array(range(0,7))*0.5 
y = 8*cos(x/1.1) + random.rand(x.shape[0]).round(2)
res4 = trapz(y,x)/3
\end{pylabcode}
$$
x = 
\begin{pylabcode} 
print(array2string(x.flatten(),separator=', ',precision=1))
\end{pylabcode} 
, \quad y = 
\begin{pylabcode} 
print(array2string(y.flatten(),separator=', ',precision=2))
\end{pylabcode}
$$
Найдите среднее значение $\hat y$ функции на интервале от 0 до 3, используя метод трапеций для вычисления интеграла. Напишите код, определяющий среднее значение функции
\fillwithdottedlines{2cm}
\answerline[$\hat y =$ \color{white}{\pylab{'{:4.2f}'.format(res4)}}] 

\question[2]Напишите код, вычисляющий интеграл:
\begin{pylabcode} 
a = int((random.rand(1,1)*10+1)[0])
b = int((random.rand(1,1)*10+1)[0])
c = int((random.rand(1,1)*10+1)[0])
d = int((random.rand(1,1)*10+1)[0])
res5 = integrate.quad(lambda x: (c+sin(x)*sin(x))/(d+a*x*x+b*(cos(x)*cos(x))), 1, 2)
\end{pylabcode}
$$
I = \int_1^2 \frac{\pylab{c} + \sin^2 x}{\pylab{d}+ \pylab{a} x^2 + \pylab{b} \cos^2 x} dx  
$$
\fillwithdottedlines{2cm}
Напишите значение интеграла:
\answerline[$I =$ \color{white}{\pylab{'{:4.2f}'.format(res5[0])}}] 

\question[3]
\begin{pylabcode} 
from scipy.optimize import linprog  
A = (random.rand(3,2)*5+2 ).astype(int) 
P = (random.rand(1,2)*40+100).astype(int)  
Z = (random.rand(3,1)*500+300).astype(int) 
res = linprog(-P, A, Z)   
res = array2string(res.x.astype(int),separator=', ') 
\end{pylabcode}
Для производства двух видов изделий (А и В) предприятие использует три вида сырья: S1, S2, S3. Для производства изделия А требуется \pylab{A[0,0]} кг сырья S1, \pylab{A[1,0]} кг сырья S2 и \pylab{A[2,0]} кг сырья S3.  Для производства изделия В требуется \pylab{A[0,1]} кг сырья S1, \pylab{A[1,1]} кг сырья S2 и \pylab{A[2,1]} кг сырья S3. Всего на предприятии имеется \pylab{Z[0,0]} кг S1, \pylab{Z[1,0]} кг S2 и \pylab{Z[2,0]} кг S3. Прибыль от реализации одного изделия А составляет \pylab{P[0,0]} руб, одного изделия В -- \pylab{P[0,1]} руб. Составить такой план продукции, при котором прибыль предприятия от реализации продукции будет максимальной.
\answerline[{\color{white}{\pylab{res}}}]
 
\question[4]
\begin{pylabcode} 
x = (random.rand(1,7)*20).astype(int)
y = (random.rand(1,7)*30).astype(int)
\end{pylabcode} 
Известны координаты семи точек на плоскости, заданных массивами координат x и y
$$
x =  
\begin{pylabcode} 
print(array2string(x.flatten(),separator=', ',precision=1)) 
\end{pylabcode}
,\quad y = 
\begin{pylabcode} 
print(array2string(y.flatten(),separator=', ',precision=1)) 
\end{pylabcode}
$$
Напишите код, который определяет и выводит на экран расстояние между двумя точками, максимально удалёнными друг от друга.
\fillwithdottedlines{7cm}
\end{questions}

\end{document}
  